\section{Conclusion}

Proof of AI is credible only when the verification claim and the consequence authority are
both narrow. Exact-integer execution gives heterogeneous hardware a deterministic
statement: a GPU may perform the original workload while a CPU validator checks sampled
matrix relations in quadratic rather than cubic time. A versioned kernel profile defines
the arithmetic, and a canonical transcript binds that arithmetic to the admitted job,
ordered graph, artifacts, outputs, availability, meter, provider, destination, and nonce.
Neither a signature, a plurality of attestations, nor cross-language byte parity proves
useful work by itself.

The reference implementations and adversarial fixtures show that this construction is
practical and that fail-closed economic gates can be tested. They are not yet a live
consensus deployment. Production activation still requires normal model-path emission,
data availability, full graph coverage, A-Chain admission and finalized-proof wiring,
nullifier and metering enforcement, and audit-gated P3Q settlement. Shadow and replay
stages can compare evidence, but must have no mint authority.

The resulting omnichain design has one source of truth. A-Chain alone validates eligible
useful work and authorizes \textsc{AI}; Z-Chain settles the finalized transition through
P3Q; Lux Quasar provides post-quantum finality; Zoo and other chains consume receipts and
may pay local rewards exactly once. This preserves the useful-work claim while avoiding a
different ``proof of whatever'' on every destination chain.
